\documentclass[11pt]{article}
\usepackage[margin=1in]{geometry}
\usepackage[T1]{fontenc}
\usepackage[utf8]{inputenc}
\usepackage{lmodern}
\usepackage{amsmath}
\usepackage{amssymb}
\usepackage{enumitem}
\usepackage{hyperref}
\hypersetup{colorlinks=true, linkcolor=black, urlcolor=blue}
\setlist[itemize]{leftmargin=1.5em}
\setlist[enumerate]{leftmargin=1.5em}

\title{Puzzle RL Theory:\\Where RL Seems To Matter Beyond SFT}
\author{Murphy}
\date{2026-04-10 23:35 UTC}

\begin{document}
\maketitle

\section*{Scope}

This note answers one narrow question: what reinforcement learning is buying over supervised fine-tuning in current LLM reasoning systems, and how to test that cleanly in a controlled puzzle setting. The target object is not ``which method gets the highest benchmark number?'' It is the mechanism split between:

\begin{itemize}
\item fixed-dataset imitation,
\item search or sampling from the base policy,
\item online data refresh from self-generated rollouts,
\item sparse outcome-reward RL,
\item and denser step-wise credit assignment.
\end{itemize}

\section*{Short answer}

\begin{enumerate}
\item \textbf{Exact theory:} the clean objective bridge is fixed-reference KL-regularized RL. In that setting, maximizing expected reward minus a KL penalty is exactly equivalent to reverse-KL projection onto a Gibbs-reweighted target policy. This is the strongest safe ``RL vs distillation'' statement because it is an identity, not an analogy.
\item \textbf{Current empirical read:} much of ordinary token-space RLVR looks more like distribution sharpening than support expansion. Several supplied papers show that strong gains can often be recovered by better sampling, online self-training, or best-of-N filtering plus distillation.
\item \textbf{Where RL still seems to matter:} when correct full trajectories are extremely rare, plain outcome-only RLVR is too weak. In that regime, denser credit assignment or more diversity-preserving exploration may add signal that fixed-trace imitation and sparse RL both miss.
\end{enumerate}

\section*{1. What is actually different about RL and SFT?}

\subsection*{Fixed-data imitation versus moving-data optimization}

SFT minimizes negative log-likelihood on a fixed dataset of prompt-response pairs. RL optimizes expected reward on trajectories sampled from the current policy. That means RL changes two things at once:

\begin{itemize}
\item the objective, because reward is not the same as token imitation;
\item the training distribution, because the rollouts come from the current policy and therefore move during training.
\end{itemize}

This second difference is easy to understate. Several of the strongest non-RL baselines in the supplied set work surprisingly well precisely because they keep the data on-policy even without using a reward model.

\subsection*{The exact KL bridge}

For prompt distribution $p(x)$, reward $r(x,y)$, reference policy $\pi_{\mathrm{ref}}(y \mid x)$, and candidate policy $\pi(y \mid x)$, consider
\[
J(\pi)=\mathbb{E}_x\left[\mathbb{E}_{y \sim \pi}[r(x,y)]-\beta\,\mathrm{KL}\!\left(\pi(\cdot \mid x)\,\|\,\pi_{\mathrm{ref}}(\cdot \mid x)\right)\right].
\]
Then the optimizer is the Gibbs-reweighted policy
\[
\pi^\star(y \mid x)\propto \pi_{\mathrm{ref}}(y \mid x)\exp(r(x,y)/\beta),
\]
and maximizing $J(\pi)$ is exactly equivalent to minimizing $\mathrm{KL}(\pi \,\|\, \pi^\star)$ up to an additive constant.

\section*{2. What the supplied papers say}

\subsection*{2.1 Evidence that a lot of the gain is not unique to RL}

\paragraph{Sampling alone can recover much of the gain.}
``Reasoning with Sampling'' argues that a training-free sampler built from the base model's own likelihoods can nearly match or beat GRPO on several reasoning tasks while preserving pass@k and diversity.\footnote{Source [4].} The key point is not that training never helps. It is that a surprisingly large amount of post-RL performance seems to live in high-likelihood regions the base model already knows how to reach.

\paragraph{Reward-free online self-training can also recover much of the gain.}
``A Model Can Help Itself'' shows that self-generated supervision alone can improve reasoning on some models without any external reward or verifier.\footnote{Source [5].} Two details are especially relevant for our project:

\begin{itemize}
\item \textbf{online refresh matters:} frozen self-generated data are much weaker than repeatedly refreshing with the latest policy;
\item \textbf{temperature decoupling matters:} low-temperature generation plus standard-temperature training works far better than coupling the two temperatures.
\end{itemize}

So one part of what people call ``RL gain'' is often just the value of staying on-policy.

\paragraph{Best-of-N positives plus cloning already go far in binary-feedback settings.}
``OREAL'' proves that, under binary outcome feedback, behavior cloning on positive trajectories from best-of-N sampling is sufficient for the KL-regularized optimal policy on the positive support.\footnote{Source [2].} Their added value comes from handling negative samples better and from adding token-level credit assignment, not from disproving the power of search-plus-distill.

\paragraph{Standard RLVR often looks like trend amplification.}
``On the Linearity of LLMs' RLVR Training'' finds strong approximate linearity in both weights and output log-probabilities over RLVR training steps.\footnote{Source [6].} Their interpretation is that RLVR often keeps amplifying directions found early, rather than continually discovering qualitatively new behavior.

\paragraph{Improved answers do not imply better traces.}
``Beyond Semantics'' is especially important for our later maze experiments.\footnote{Source [3].} In a controlled puzzle-style setting, the paper shows:

\begin{itemize}
\item trace validity is not a reliable predictor of final correctness;
\item corrupted traces can train almost as well as correct ones, and can generalize better;
\item GRPO can improve answer quality without improving trace validity.
\end{itemize}

\subsection*{2.2 Where RL-like learning still seems to add something real}

\paragraph{When pass@k is effectively zero, sparse RL is too weak.}
``Supervised Reinforcement Learning'' targets the regime where the model almost never samples a fully correct trajectory, so outcome-only RLVR has no useful signal and plain SFT overfits long demonstrations.\footnote{Source [7].} Their dense step-wise similarity reward over expert actions is best read as evidence that \emph{credit assignment} can be the real bottleneck.

\paragraph{Diversity-preserving exploration may matter.}
``Beyond Mode Elicitation'' directly targets diversity collapse.\footnote{Source [8].} Their main claim is stronger than a pass@1 win: latent-space exploration improves both pass@1 and pass@k relative to discrete RL baselines. If that result generalizes, it is one of the better candidates for ``RL is doing more than just sharpening one dominant mode.''

\section*{3. Best current synthesis}

\subsection*{What seems safely established}

\begin{itemize}
\item The exact theoretical bridge is fixed-reference KL-RL, not a blanket equivalence between all RL and all SFT.
\item A large fraction of current token-space RLVR gains can be explained by sampling, verifier filtering, and online data refresh over behaviors already latent in the base model.
\item Plain final-answer gains do not by themselves show better process quality or support expansion.
\item Dense or structured credit assignment matters when correct full trajectories are too rare for outcome-only RL to learn from.
\item Diversity must be treated as part of the scientific object, not a side metric.
\end{itemize}

\subsection*{What seems plausible but not yet proved}

\begin{itemize}
\item Some diversity-preserving RL variants may truly enlarge the set of successful reasoning trajectories rather than only reweighting one mode.
\item Latent-space exploration may be a better optimizer family than discrete token RL when the main failure mode is entropy collapse.
\item In long-horizon agentic settings, on-policy learning may matter more than it does in single-turn math or code.
\end{itemize}

\subsection*{What is still open}

\begin{itemize}
\item How often does RL find successful trajectories that large-k base sampling could not have found?
\item When RL helps, is the key ingredient reward optimization, on-policy refresh, verifier filtering, or dense credit assignment?
\item How much of any RL gain compresses back into an ordinary supervised dataset?
\item What should count as genuine ``distribution expansion'': lower base-policy likelihood, new algorithm family, stronger OOD transfer, or all three?
\end{itemize}

\section*{4. What this means for the maze repo}

The existing repo already has the right seed object:

\begin{itemize}
\item graph generation and shortest-path verification,
\item supervised pretrain and midtrain stages,
\item path validity and path optimality metrics.
\end{itemize}

What it does \emph{not} yet have is the mechanism surface we need. The first implementation order should be:

\begin{enumerate}
\item add multi-sample evaluation with pass@k, diversity logging, and base-logprob logging;
\item add best-of-N filtering plus behavior cloning;
\item add reward-free online self-training;
\item only then add sparse binary-outcome RLVR;
\item add a step-wise dense-credit baseline after that;
\item treat diversity-preserving RL as a later-stage extension.
\end{enumerate}

\section*{5. Falsifiable puzzle hypotheses}

\begin{enumerate}
\item \textbf{Mostly-sharpening hypothesis:} on the current maze family, best-of-N plus cloning and reward-free online self-training will recover most of the pass@1 gain that sparse RLVR gets.
\item \textbf{Credit-assignment hypothesis:} on harder maze regimes where base pass@k is near zero, dense step-wise feedback will beat sparse outcome-only RLVR.
\item \textbf{Diversity hypothesis:} standard discrete RLVR will tend to improve pass@1 at the cost of pass@k, while any truly better exploration method must avoid that trade.
\item \textbf{Compression hypothesis:} if a gain is mostly due to better trajectory discovery rather than optimizer-specific magic, much of it should survive distillation back into a supervised dataset.
\end{enumerate}

\section*{Bottom line}

The strongest current reading is not ``RL does nothing'' and not ``RL creates brand-new reasoning by default.'' It is more specific:

\begin{quote}
Current LLM RL often behaves like on-policy distribution shaping over capabilities that are already latent in the base model. The main exceptions appear when better data refresh, better verifier-filtered search, or better credit assignment materially change the learning problem. The maze project should be built to separate those ingredients one by one.
\end{quote}

\section*{Sources}

\begin{enumerate}
\item Internal bridge note: \texttt{projects/llm-rl-kl-duality/docs/theory-bridge.md}
\item \href{https://arxiv.org/abs/2502.06781}{arXiv:2502.06781} -- \emph{Exploring the Limit of Outcome Reward for Learning Mathematical Reasoning}
\item \href{https://arxiv.org/abs/2505.13775}{arXiv:2505.13775} -- \emph{Beyond Semantics: The Unreasonable Effectiveness of Reasonless Intermediate Tokens}
\item \href{https://arxiv.org/abs/2510.14901}{arXiv:2510.14901} -- \emph{Reasoning with Sampling: Your Base Model is Smarter Than You Think}
\item \href{https://arxiv.org/abs/2510.18814}{arXiv:2510.18814} -- \emph{A Model Can Help Itself: Reward-Free Self-Training for LLM Reasoning}
\item \href{https://arxiv.org/abs/2601.04537}{arXiv:2601.04537} -- \emph{Not All Steps are Informative: On the Linearity of LLMs' RLVR Training}
\item \href{https://arxiv.org/abs/2510.25992}{arXiv:2510.25992} -- \emph{Supervised Reinforcement Learning: From Expert Trajectories to Step-wise Reasoning}
\item \href{https://arxiv.org/abs/2602.01705}{arXiv:2602.01705} -- \emph{Beyond Mode Elicitation: Diversity-Preserving Reinforcement Learning via Latent Diffusion Reasoner}
\item \href{https://arxiv.org/abs/2507.10616}{arXiv:2507.10616} -- \emph{Scalpel vs. Hammer: GRPO Amplifies Existing Capabilities, SFT Replaces Them}
\end{enumerate}

\end{document}
